\documentclass[12pt]{article}
\usepackage[utf8]{inputenc}
\usepackage{amsmath, amssymb, amsthm}
\usepackage{geometry}
\usepackage{xcolor}
\usepackage[most]{tcolorbox}
\usepackage{enumitem}
\usepackage{fancyhdr}

% Page setup
\geometry{margin=1in}
\setlength{\headheight}{14.5pt}
\pagestyle{fancy}
\fancyhf{}
\rhead{AMC12/COMC Problem Analysis}
\lhead{\today}

% Color definitions
\definecolor{problemconfig}{RGB}{230, 242, 255}    % Light blue
\definecolor{strategic}{RGB}{255, 230, 242}       % Light pink
\definecolor{tactical}{RGB}{242, 255, 230}        % Light green
\definecolor{techniques}{RGB}{255, 242, 230}      % Light yellow
\definecolor{verification}{RGB}{255, 230, 230}    % Light red
\definecolor{solution}{RGB}{230, 255, 255}        % Light cyan
\definecolor{competition}{RGB}{242, 242, 255}     % Light purple
\definecolor{gold}{RGB}{218, 165, 32}

% Box styles
\newtcolorbox{problemconfigbox}{
    colback=problemconfig,
    colframe=blue,
    title=Problem Configuration Analysis,
    fonttitle=\bfseries,
    arc=3pt,
    boxrule=1pt,
    breakable
}

\newtcolorbox{strategicbox}{
    colback=strategic,
    colframe=purple,
    title=Strategic Framework Application,
    fonttitle=\bfseries,
    arc=3pt,
    boxrule=1pt,
    breakable
}

\newtcolorbox{tacticalbox}{
    colback=tactical,
    colframe=green,
    title=Tactical Methodology Selection,
    fonttitle=\bfseries,
    arc=3pt,
    boxrule=1pt,
    breakable
}

\newtcolorbox{techniquesbox}{
    colback=techniques,
    colframe=orange,
    title=Conversion Techniques Application,
    fonttitle=\bfseries,
    arc=3pt,
    boxrule=1pt,
    breakable
}

\newtcolorbox{verificationbox}{
    colback=verification,
    colframe=red,
    title=Verification Strategy Design,
    fonttitle=\bfseries,
    arc=3pt,
    boxrule=1pt,
    breakable
}

\newtcolorbox{solutionbox}{
    colback=solution,
    colframe=cyan,
    title=Solution Roadmap,
    fonttitle=\bfseries,
    arc=3pt,
    boxrule=1pt,
    breakable
}

\newtcolorbox{competitionbox}{
    colback=competition,
    colframe=violet,
    title=Competition Strategy Integration,
    fonttitle=\bfseries,
    arc=3pt,
    boxrule=1pt,
    breakable
}

\newtcolorbox{keyinsight}{
    colback=yellow!10,
    colframe=gold,
    title=\textcolor{gold}{$\star$ Key Insight},
    fonttitle=\bfseries,
    arc=3pt,
    boxrule=2pt,
    breakable
}

\newtcolorbox{warning}{
    colback=red!5,
    colframe=red!70!black,
    title=\textcolor{red!70!black}{$\triangle$ Warning},
    fonttitle=\bfseries,
    arc=3pt,
    boxrule=2pt,
    breakable
}

\newtcolorbox{protip}{
    colback=green!5,
    colframe=green!70!black,
    title=\textcolor{green!70!black}{$\checkmark$ Pro Tip},
    fonttitle=\bfseries,
    arc=3pt,
    boxrule=2pt,
    breakable
}

\begin{document}

\title{\textbf{Strategic Analysis: AMC 12B 2024 Problem 23}}
\author{Comprehensive Problem-Solving Framework}
\date{\today}
\maketitle

\section*{Problem Statement}

A right pyramid has regular octagon $ABCDEFGH$ with side length $1$ as its base and apex $V$. Segments $\overline{AV}$ and $\overline{DV}$ are perpendicular. What is the square of the height of the pyramid?

$$\textbf{(A) } 1 \qquad \textbf{(B) } \frac{1+\sqrt{2}}{2} \qquad \textbf{(C) } \sqrt{2} \qquad \textbf{(D) } \frac{3}{2} \qquad \textbf{(E) } \frac{2+\sqrt{2}}{3}$$

\vspace{1cm}

\begin{problemconfigbox}
\subsection*{1. Problem Configuration Analysis}

\subsubsection*{A. Problem Classification}

\begin{itemize}[leftmargin=*]
    \item \textbf{Problem Type}: To Find (computational with geometric constraint)
    \item \textbf{Difficulty Band}: Late-band (Problem 23/25)
    \item \textbf{Estimated Time}: 6-8 minutes
    \item \textbf{Competition Context}: AMC 12B, multiple choice format
    \item \textbf{Mathematical Domain}: 3D Geometry (solid geometry, pyramid) with Coordinate Geometry and Trigonometry
    \item \textbf{Cross-Domain Nature}: This is a sophisticated cross-domain problem requiring:
    \begin{itemize}
        \item \textbf{3D Geometry}: Pyramid structure, perpendicularity in space
        \item \textbf{2D Geometry}: Regular octagon properties, symmetry
        \item \textbf{Coordinate Geometry}: Possible coordinate approach
        \item \textbf{Vectors}: Dot product for perpendicularity
        \item \textbf{Algebra}: System solving and manipulation
        \item \textbf{Trigonometry}: Law of Cosines in octagon
    \end{itemize}
\end{itemize}

\subsubsection*{B. Problem Structure Identification}

\textbf{Given Information}:
\begin{itemize}
    \item Regular octagon $ABCDEFGH$ with side length $s = 1$
    \item Right pyramid: apex $V$ directly above center of octagon
    \item Constraint: $\overline{AV} \perp \overline{DV}$ (perpendicular segments)
    \item Vertices labeled consecutively around octagon
    \item Multiple choice answers provided (all involving $\sqrt{2}$)
\end{itemize}

\textbf{Constraints}:
\begin{itemize}
    \item Regular octagon: all sides equal length, all interior angles equal ($135^\circ$)
    \item Right pyramid: base is horizontal, apex directly above center
    \item By symmetry: $V$ is equidistant from all vertices of octagon
    \item Perpendicularity constraint: $\overrightarrow{AV} \cdot \overrightarrow{DV} = 0$
    \item $A$ and $D$ are separated by 3 edges (vertices $A, B, C, D$)
\end{itemize}

\textbf{Objective}:
\begin{itemize}
    \item Find $h^2$ where $h$ is the height of the pyramid (distance from center $O$ to apex $V$)
    \item Note: Problem asks for the \textit{square} of the height, not the height itself
\end{itemize}

\textbf{Hidden Assumptions}:
\begin{itemize}
    \item Center of octagon is denoted $O$ (not explicitly stated)
    \item Octagon lies in horizontal plane
    \item The perpendicularity condition uniquely determines the height
    \item All answer choices are algebraic numbers involving $\sqrt{2}$
\end{itemize}

\textbf{Answer Choice Leverage}:
\begin{itemize}
    \item All involve $\sqrt{2}$ (suggests octagon geometry is key)
    \item Answer (B) has form $\frac{1+\sqrt{2}}{2}$ (simplest numerator)
    \item Answers range from 1 to about 1.5, giving bounds
    \item Can check dimensionality: must have units of length$^2$
\end{itemize}

\subsubsection*{C. Strategic Orientation}

\textbf{Hypothesis}: The perpendicularity constraint $AV \perp DV$ determines a unique height $h$ for the pyramid.

\textbf{Conclusion}: Compute $h^2$ using geometric relationships.

\textbf{Notation}:
\begin{itemize}
    \item $O$ = center of octagon (origin)
    \item $h$ = height of pyramid = $|OV|$
    \item $r$ = radius of octagon = distance from center to vertex
    \item $A, B, C, D, \ldots, H$ = vertices of octagon
    \item $V$ = apex of pyramid
\end{itemize}

\textbf{Intuitive Approach}:
\begin{itemize}
    \item Find distance $|AD|$ using octagon geometry
    \item Use perpendicularity: $\triangle AVD$ is right-angled at $V$
    \item Apply Pythagorean theorem in 3D
    \item Relate $|AV|$, $|AD|$, and $h$ through geometric constraints
\end{itemize}

\textbf{Cross-Domain Potential}:
\begin{itemize}
    \item \textbf{Pure Geometry}: Use symmetry and special triangles
    \item \textbf{Coordinate Geometry}: Place octagon in $xy$-plane, use 3D coordinates
    \item \textbf{Vector Approach}: Use dot product for perpendicularity
    \item \textbf{Analytical Geometry}: Compute all distances algebraically
\end{itemize}

\end{problemconfigbox}

\begin{strategicbox}
\subsection*{2. Strategic Framework Application}

\subsubsection*{A. Psychological Strategies}

\textbf{Mental Toughness}:
\begin{itemize}
    \item This is Problem 23 - expect sophisticated 3D geometry
    \item Visualization is key but may be challenging
    \item Multiple approaches available - be flexible
    \item Don't panic if first approach seems complex
\end{itemize}

\textbf{Creativity Cultivation}:
\begin{itemize}
    \item Consider: Can I reduce this to 2D problem?
    \item Think: What special properties does a regular octagon have?
    \item Explore: How does perpendicularity in 3D constrain the height?
    \item Question: Can symmetry simplify the problem?
\end{itemize}

\textbf{Iterative Refinement}:
\begin{itemize}
    \item First: Understand octagon structure (inscribed in circle)
    \item Then: Find key distances ($|AD|$, radius $r$)
    \item Next: Use perpendicularity to set up equation
    \item Finally: Solve for $h^2$
\end{itemize}

\subsubsection*{B. Investigation Strategies}

\textbf{Startup Orientation}:
\begin{itemize}
    \item Draw diagram of octagon with labeled vertices
    \item Visualize pyramid in 3D
    \item Identify that $A$ and $D$ are 3 vertices apart
\end{itemize}

\textbf{Get Your Hands Dirty}:
\begin{itemize}
    \item Compute: What is $|AD|$ for regular octagon with side 1?
    \item Compute: What is radius $r$ of circumscribed circle?
    \item Explore: If $AV \perp DV$, what does this say about $\triangle AVD$?
\end{itemize}

\textbf{Penultimate Step}:
\begin{itemize}
    \item Final step: Extract $h^2$ from equation
    \item Penultimate step: Set up Pythagorean relation in pyramid
    \item Working backward: Need $|AV|$ in terms of known quantities
\end{itemize}

\textbf{Wishful Thinking}:
\begin{itemize}
    \item Wish: "I want $\triangle AVD$ to be a nice triangle"
    \item Reality: Since $AV \perp DV$ and by symmetry $|AV| = |DV|$, it's 45-45-90!
    \item Wish: "I want coordinates to make calculations easier"
    \item Strategy: Use coordinate system aligned with octagon
\end{itemize}

\textbf{Make It Easier}:
\begin{itemize}
    \item Simplify: Use symmetry ($|AV| = |DV|$ by pyramid symmetry)
    \item Reduce dimension: Work with projections and right triangles
    \item Scale: Could scale octagon by convenient factor if helpful
\end{itemize}

\textbf{Conjecture via Small Cases}:
\begin{itemize}
    \item Test extreme: If $h = 0$ (flat), would $AV \perp DV$? No.
    \item Test extreme: If $h$ very large, would $AV \perp DV$? No.
    \item Conclusion: There's a specific $h$ satisfying the constraint
\end{itemize}

\textbf{Visualization}:
\begin{itemize}
    \item Draw octagon from above
    \item Draw pyramid from side view
    \item Visualize the constraint: two segments from apex meeting at right angle
\end{itemize}

\subsubsection*{C. Late-Band Strategic Playbook}

\textbf{Answer-Choice Leverage}:
\begin{itemize}
    \item All answers contain $\sqrt{2}$: octagon geometry crucial
    \item Can eliminate by dimensional analysis
    \item Can check special cases or limits
    \item Range check: $h^2$ should be reasonable (not too large or small)
\end{itemize}

\textbf{Make It Easier}:
\begin{itemize}
    \item Focus on key triangle $AVD$ rather than entire pyramid
    \item Use symmetry to avoid redundant calculations
    \item Find one key distance that unlocks everything
\end{itemize}

\textbf{Penultimate Step}:
\begin{itemize}
    \item Work backwards: Given $h^2 = \frac{1+\sqrt{2}}{2}$, verify it works
    \item This can be faster than deriving from scratch
\end{itemize}

\textbf{Wishful Thinking}:
\begin{itemize}
    \item Assume $\triangle AVD$ is 45-45-90 (turns out to be true!)
    \item This immediately gives $|AV| = |DV| = \frac{|AD|}{\sqrt{2}}$
\end{itemize}

\textbf{Conjecture via Coordinates}:
\begin{itemize}
    \item Set up coordinates and compute dot product
    \item This is methodical and reliable for 3D problems
\end{itemize}

\end{strategicbox}

\begin{tacticalbox}
\subsection*{3. Tactical Methodology Selection}

\subsubsection*{A. Core Mathematical Tactics}

\textbf{Primary Tactics}:
\begin{enumerate}
    \item \textbf{Regular Octagon Properties}:
    \begin{itemize}
        \item Interior angle: $\frac{(8-2) \cdot 180^\circ}{8} = 135^\circ$
        \item Central angle: $\frac{360^\circ}{8} = 45^\circ$
        \item Can decompose into 8 congruent isosceles triangles
        \item Inscribed in circle with radius $r$
    \end{itemize}
    
    \item \textbf{3D Pythagorean Theorem}:
    $$|AV|^2 = |AO|^2 + |OV|^2 = r^2 + h^2$$
    
    \item \textbf{Perpendicularity in 3D}:
    For vectors $\vec{u}$ and $\vec{v}$: $\vec{u} \perp \vec{v} \Leftrightarrow \vec{u} \cdot \vec{v} = 0$
    
    \item \textbf{Symmetry Exploitation}:
    By pyramid symmetry: $|AV| = |BV| = \cdots = |HV|$ (all equal)
    
    \item \textbf{Special Triangle Recognition}:
    If $|AV| = |DV|$ and $AV \perp DV$, then $\triangle AVD$ is 45-45-90
\end{enumerate}

\textbf{Supporting Tactics}:
\begin{itemize}
    \item \textbf{Law of Cosines}: For finding $|AD|$ in octagon
    \item \textbf{Coordinate Geometry}: Setting up $(x, y, z)$ coordinates
    \item \textbf{Distance Formula}: Computing distances in 3D
    \item \textbf{Isosceles Triangle Properties}: Utilizing symmetry
\end{itemize}

\subsubsection*{B. Crossover Techniques}

\begin{itemize}
    \item \textbf{3D Geometry $\to$ 2D Projections}: Reduce to octagon plane
    \item \textbf{Geometry $\to$ Algebra}: Convert to coordinate system
    \item \textbf{Vectors}: Use dot product for perpendicularity
    \item \textbf{Special Angles}: Leverage $45^\circ$ angles in octagon
    \item \textbf{Symmetry $\to$ Simplification}: Reduce problem complexity
\end{itemize}

\end{tacticalbox}

\begin{techniquesbox}
\subsection*{4. Conversion Techniques Application}

\subsubsection*{A. Recognition Index}

\textbf{Triggered Techniques}:
\begin{itemize}
    \item \textbf{Regular Polygon}: Use central angle = $360^\circ/n$
    \item \textbf{Right Pyramid}: Apex directly above center
    \item \textbf{Perpendicularity}: Use dot product or triangle properties
    \item \textbf{Symmetry}: All vertices equidistant from apex
    \item \textbf{45-45-90 Triangle}: Special right triangle with ratio $1:1:\sqrt{2}$
\end{itemize}

\subsubsection*{B. Technique Selection Priority}

\textbf{Primary Path} (Recommended - Geometric):
\begin{enumerate}
    \item Find $|AD|$ using octagon geometry (Law of Cosines or direct computation)
    \item Recognize $\triangle AVD$ is isosceles ($|AV| = |DV|$ by symmetry)
    \item Since $AV \perp DV$, conclude $\triangle AVD$ is 45-45-90
    \item Calculate $|AV| = \frac{|AD|}{\sqrt{2}}$
    \item Use 3D Pythagorean: $|AV|^2 = r^2 + h^2$
    \item Solve for $h^2$
\end{enumerate}

\textbf{Alternative Path 1} (Coordinate/Vector Approach):
\begin{enumerate}
    \item Set up coordinates with octagon in $xy$-plane
    \item Place $A$ and $D$ at calculated positions
    \item Apex $V$ at $(0, 0, h)$
    \item Compute $\vec{AV}$ and $\vec{DV}$
    \item Use $\vec{AV} \cdot \vec{DV} = 0$ to solve for $h$
\end{enumerate}

\textbf{Alternative Path 2} (Scaling Approach):
\begin{enumerate}
    \item Scale octagon by factor $k$ for easier arithmetic
    \item Perform all calculations with scaled octagon
    \item Scale answer back by factor $k^2$
\end{enumerate}

\subsubsection*{C. Integration Strategy}

\textbf{Key Integration}: Combining 2D octagon geometry with 3D pyramid structure

\textbf{Conversion Sequence}:
$$\text{Regular Octagon} \xrightarrow{\text{Geometry}} |AD|, r \xrightarrow{\text{Symmetry}} |AV| = |DV|$$
$$\xrightarrow{\text{Perpendicularity}} 45\text{-}45\text{-}90 \xrightarrow{\text{3D Pythagorean}} h^2$$

\textbf{Critical Insights}:
\begin{itemize}
    \item The regular octagon has $A$ and $D$ separated by 3 edges
    \item Central angle subtended by arc $AD$ is $3 \times 45^\circ = 135^\circ$
    \item By symmetry of right pyramid, all slant edges equal
    \item Perpendicularity + equality $\Rightarrow$ 45-45-90 triangle
\end{itemize}

\subsubsection*{D. Conflict Resolution}

\textbf{Potential Conflicts}:
\begin{itemize}
    \item 3D visualization vs. algebraic calculation
    \item Multiple coordinate system choices
    \item Different formulas for octagon radius
\end{itemize}

\textbf{Resolution Strategy}:
\begin{itemize}
    \item Use most natural coordinate system (octagon center at origin)
    \item Verify calculations with dimensional analysis
    \item Check answer against bounds and symmetry
\end{itemize}

\end{techniquesbox}

\begin{verificationbox}
\subsection*{5. Verification Strategy Design}

\subsubsection*{A. Verification Level Selection}

\textbf{Recommended Level}: Level 4-5 (3-5 minutes)

\textbf{Rationale}:
\begin{itemize}
    \item Problem 23 is high-value (late-band)
    \item 3D geometry prone to visualization errors
    \item Multiple calculations that need checking
    \item Perpendicularity constraint must be verified
\end{itemize}

\subsubsection*{B. Specialized Verification Methods}

\textbf{For Final Answer $h^2 = \frac{1+\sqrt{2}}{2}$}:

\begin{enumerate}
    \item \textbf{Verify Octagon Calculations}:
    \begin{itemize}
        \item Check: $|AD| = 1 + \sqrt{2}$ (distance between vertices 3 apart)
        \item Check: Radius $r^2 = \frac{2+\sqrt{2}}{2}$ (from Law of Cosines)
    \end{itemize}
    
    \item \textbf{Verify Triangle $AVD$}:
    \begin{itemize}
        \item $|AV| = |DV| = \frac{1+\sqrt{2}}{\sqrt{2}} = \frac{\sqrt{2}+2}{2}$ (by 45-45-90)
        \item Check: $|AV|^2 = \frac{(2+\sqrt{2})^2}{4} = \frac{6+4\sqrt{2}}{4} = \frac{3+2\sqrt{2}}{2}$
    \end{itemize}
    
    \item \textbf{Verify 3D Pythagorean}:
    \begin{align*}
    h^2 &= |AV|^2 - r^2 \\
    &= \frac{3+2\sqrt{2}}{2} - \frac{2+\sqrt{2}}{2} \\
    &= \frac{1+\sqrt{2}}{2} \quad \checkmark
    \end{align*}
    
    \item \textbf{Verify Perpendicularity} (if using vectors):
    Check that $\vec{AV} \cdot \vec{DV} = 0$ with computed $h$
    
    \item \textbf{Dimensional Analysis}:
    $h^2$ has dimensions of length$^2$ \checkmark
    
    \item \textbf{Reasonableness Check}:
    $h^2 = \frac{1+\sqrt{2}}{2} \approx \frac{2.414}{2} \approx 1.207$
    
    So $h \approx 1.1$, which is reasonable for octagon with side 1
    
    \item \textbf{Answer Choice Match}:
    $h^2 = \frac{1+\sqrt{2}}{2} = \text{Choice (B)}$ \quad \checkmark
\end{enumerate}

\subsubsection*{C. Confidence Calibration}

\textbf{Confidence Assessment}: 90-95\%

\textbf{Reasons for High Confidence}:
\begin{itemize}
    \item Geometric reasoning is sound (45-45-90 triangle)
    \item All intermediate calculations check out
    \item Dimensional analysis correct
    \item Answer is among choices
    \item Multiple solution methods converge to same answer
\end{itemize}

\textbf{Residual Uncertainty}:
\begin{itemize}
    \item 3D visualization could have subtle errors (5\%)
    \item Arithmetic errors in radical manipulation (5\%)
\end{itemize}

\end{verificationbox}

\begin{keyinsight}
\textbf{Central Insight}: By symmetry of the right pyramid, $|AV| = |DV|$. Combined with the perpendicularity constraint $AV \perp DV$, this makes $\triangle AVD$ a 45-45-90 right triangle. This elegant geometric fact allows us to express $|AV|$ in terms of $|AD|$, which can be computed from octagon geometry. The 3D Pythagorean theorem then directly gives $h^2$.
\end{keyinsight}

\begin{solutionbox}
\subsection*{6. Solution Roadmap}

\subsubsection*{Step-by-Step Solution}

\textbf{Step 1: Find Distance $|AD|$ in Regular Octagon}

Consider regular octagon $ABCDEFGH$ with side length 1. Vertices $A$ and $D$ are separated by 3 edges.

The octagon can be inscribed in a circle. The central angle subtended by arc from $A$ to $D$ (going through $B$ and $C$) is:
$$\theta = 3 \times \frac{360^\circ}{8} = 3 \times 45^\circ = 135^\circ$$

Let $O$ be the center and $r$ be the circumradius. First, find $r$ using one side.

For side $AB$, the central angle is $45^\circ$. Using the Law of Cosines in $\triangle AOB$:
$$|AB|^2 = r^2 + r^2 - 2r^2\cos(45^\circ) = 2r^2(1 - \cos 45^\circ) = 2r^2\left(1 - \frac{\sqrt{2}}{2}\right)$$

Since $|AB| = 1$:
$$1 = 2r^2\left(1 - \frac{\sqrt{2}}{2}\right) = 2r^2 \cdot \frac{2-\sqrt{2}}{2} = r^2(2-\sqrt{2})$$

Therefore:
$$r^2 = \frac{1}{2-\sqrt{2}} = \frac{1}{2-\sqrt{2}} \cdot \frac{2+\sqrt{2}}{2+\sqrt{2}} = \frac{2+\sqrt{2}}{4-2} = \frac{2+\sqrt{2}}{2}$$

Now find $|AD|$ using Law of Cosines in $\triangle AOD$ with central angle $135^\circ$:
\begin{align*}
|AD|^2 &= r^2 + r^2 - 2r^2\cos(135^\circ) \\
&= 2r^2(1 - \cos 135^\circ) \\
&= 2r^2\left(1 + \frac{\sqrt{2}}{2}\right) \\
&= 2r^2 \cdot \frac{2+\sqrt{2}}{2} \\
&= r^2(2+\sqrt{2})
\end{align*}

Substituting $r^2 = \frac{2+\sqrt{2}}{2}$:
$$|AD|^2 = \frac{2+\sqrt{2}}{2} \cdot (2+\sqrt{2}) = \frac{(2+\sqrt{2})^2}{2} = \frac{4 + 4\sqrt{2} + 2}{2} = \frac{6+4\sqrt{2}}{2} = 3 + 2\sqrt{2}$$

Therefore:
$$|AD| = \sqrt{3+2\sqrt{2}} = \sqrt{(\sqrt{2}+1)^2} = \sqrt{2} + 1$$

\textbf{Step 2: Analyze Triangle $AVD$}

Since the pyramid is right (apex $V$ directly above center $O$), by symmetry:
$$|AV| = |BV| = |CV| = \cdots = |HV|$$

In particular, $|AV| = |DV|$.

Given constraint: $\overline{AV} \perp \overline{DV}$

Since $\triangle AVD$ is isosceles with $|AV| = |DV|$ and has a right angle at $V$, it must be a \textbf{45-45-90 triangle}.

In a 45-45-90 triangle with hypotenuse $c$, the legs have length $\frac{c}{\sqrt{2}}$.

Therefore:
$$|AV| = |DV| = \frac{|AD|}{\sqrt{2}} = \frac{\sqrt{2}+1}{\sqrt{2}} = \frac{(\sqrt{2}+1) \cdot \sqrt{2}}{2} = \frac{2 + \sqrt{2}}{2}$$

\textbf{Step 3: Apply 3D Pythagorean Theorem}

Consider the right triangle $\triangle AOV$ where:
\begin{itemize}
    \item $O$ is the center of the octagon (on the base)
    \item $A$ is a vertex of the octagon
    \item $V$ is the apex of the pyramid
    \item The right angle is at $O$ (since the pyramid is right)
\end{itemize}

By the Pythagorean theorem in 3D:
$$|AV|^2 = |AO|^2 + |OV|^2 = r^2 + h^2$$

where $h = |OV|$ is the height of the pyramid.

Solving for $h^2$:
$$h^2 = |AV|^2 - r^2$$

Calculate $|AV|^2$:
$$|AV|^2 = \left(\frac{2+\sqrt{2}}{2}\right)^2 = \frac{(2+\sqrt{2})^2}{4} = \frac{4 + 4\sqrt{2} + 2}{4} = \frac{6+4\sqrt{2}}{4} = \frac{3+2\sqrt{2}}{2}$$

Therefore:
\begin{align*}
h^2 &= \frac{3+2\sqrt{2}}{2} - \frac{2+\sqrt{2}}{2} \\
&= \frac{(3+2\sqrt{2}) - (2+\sqrt{2})}{2} \\
&= \frac{1 + \sqrt{2}}{2}
\end{align*}

\textbf{Answer}: $h^2 = \boxed{\frac{1+\sqrt{2}}{2}}$, which is choice $\boxed{\textbf{(B)}}$.

\subsubsection*{Alternative Solution: Vector Approach}

Set up coordinates with octagon center at origin, octagon in $xy$-plane.

Place vertices at angles $0^\circ, 45^\circ, 90^\circ, \ldots, 315^\circ$ from positive $x$-axis.

With circumradius $r$ where $r^2 = \frac{2+\sqrt{2}}{2}$:

Let $A$ be at angle $0^\circ$: $A = (r, 0, 0)$

Let $D$ be at angle $135^\circ$: $D = (-r\frac{\sqrt{2}}{2}, r\frac{\sqrt{2}}{2}, 0)$

Apex: $V = (0, 0, h)$

Vectors:
\begin{align*}
\vec{AV} &= V - A = (-r, 0, h) \\
\vec{DV} &= V - D = \left(r\frac{\sqrt{2}}{2}, -r\frac{\sqrt{2}}{2}, h\right)
\end{align*}

Perpendicularity condition:
\begin{align*}
\vec{AV} \cdot \vec{DV} &= 0 \\
(-r)\left(r\frac{\sqrt{2}}{2}\right) + (0)\left(-r\frac{\sqrt{2}}{2}\right) + h \cdot h &= 0 \\
-\frac{r^2\sqrt{2}}{2} + h^2 &= 0 \\
h^2 &= \frac{r^2\sqrt{2}}{2}
\end{align*}

Substituting $r^2 = \frac{2+\sqrt{2}}{2}$:
$$h^2 = \frac{\sqrt{2}}{2} \cdot \frac{2+\sqrt{2}}{2} = \frac{\sqrt{2}(2+\sqrt{2})}{4} = \frac{2\sqrt{2} + 2}{4} = \frac{2(\sqrt{2}+1)}{4} = \frac{\sqrt{2}+1}{2}$$

Therefore $h^2 = \boxed{\frac{1+\sqrt{2}}{2}}$. \quad \checkmark

\end{solutionbox}

\begin{warning}
\textbf{Common Pitfalls}:
\begin{itemize}
    \item \textbf{Confusing vertices}: $A$ and $D$ are 3 edges apart, not adjacent or opposite
    \item \textbf{Wrong central angle}: Arc $AD$ subtends $135^\circ$, not $90^\circ$ or $180^\circ$
    \item \textbf{Forgetting symmetry}: Must use $|AV| = |DV|$ from pyramid symmetry
    \item \textbf{Arithmetic with radicals}: Carefully expand $(2+\sqrt{2})^2 = 6 + 4\sqrt{2}$, not $6 + 2\sqrt{2}$
    \item \textbf{Coordinate placement}: In vector approach, ensure vertices are correctly positioned
    \item \textbf{3D Pythagorean}: Remember $|AV|^2 = r^2 + h^2$, not $|AD|^2$
\end{itemize}
\end{warning}

\begin{protip}
\textbf{Time-Saving Strategies}:
\begin{itemize}
    \item \textbf{Recognize 45-45-90}: As soon as you see $|AV| = |DV|$ and $AV \perp DV$, conclude it's 45-45-90 immediately
    \item \textbf{Use formula}: For regular $n$-gon with side $s$ inscribed in circle: $r = \frac{s}{2\sin(\pi/n)}$
    \item \textbf{Check answer form}: All choices have $\sqrt{2}$, confirming octagon geometry is key
    \item \textbf{Dimensional analysis}: Quick check that $h^2$ has correct dimensions
    \item \textbf{Symmetry shortcuts}: Leverage pyramid symmetry to avoid computing all vertices
    \item \textbf{Simplify radicals}: Recognize $\sqrt{3+2\sqrt{2}} = 1 + \sqrt{2}$ by pattern $(a+b)^2 = a^2 + 2ab + b^2$
\end{itemize}
\end{protip}

\begin{competitionbox}
\subsection*{7. Competition Strategy Integration}

\subsubsection*{A. Time Management}

\textbf{Time Budget}: 6-8 minutes for this P23 problem

\textbf{Allocation}:
\begin{itemize}
    \item Problem reading and visualization: 1 minute
    \item Find octagon distances ($|AD|$, $r$): 2-3 minutes
    \item Recognize 45-45-90 triangle: 30 seconds
    \item Apply 3D Pythagorean and solve: 2 minutes
    \item Verification: 1-2 minutes
\end{itemize}

\textbf{Decision Points}:
\begin{itemize}
    \item If octagon calculation stalls after 3 minutes, try coordinate/vector approach
    \item If visualization is difficult, commit to coordinate method from start
    \item If radical arithmetic is messy, carefully re-compute each step
\end{itemize}

\subsubsection*{B. Prioritization Logic}

\textbf{Why Attempt This Problem}:
\begin{itemize}
    \item Clear geometric structure (regular octagon, right pyramid)
    \item Key insight (45-45-90 triangle) makes it tractable
    \item Multiple approaches available
    \item Worth 6 points on AMC 12
    \item If comfortable with 3D geometry, this is accessible
\end{itemize}

\textbf{When to Skip}:
\begin{itemize}
    \item If uncomfortable with 3D visualization
    \item If weak in working with radicals
    \item If running very low on time (<8 minutes remaining)
    \item If easier Problem 24 or 25 appears more tractable
\end{itemize}

\subsubsection*{C. Risk Management}

\textbf{Confidence Indicators}:
\begin{itemize}
    \item Successfully found $|AD| = 1 + \sqrt{2}$ (70\% confidence)
    \item Recognized 45-45-90 triangle structure (85\% confidence)
    \item Computed $r^2$ correctly (80\% confidence)
    \item Applied 3D Pythagorean correctly (90\% confidence)
    \item Final answer matches choice (95\% confidence)
\end{itemize}

\textbf{Fallback Strategies}:
\begin{itemize}
    \item If geometric approach unclear, switch to coordinate/vector method
    \item If algebra gets messy, try numerical approximation to check answer
    \item Can verify answer by working backwards from choices
\end{itemize}

\subsubsection*{D. Strategic Abandonment}

\textbf{Abandon If}:
\begin{itemize}
    \item After 5 minutes, no progress on finding key distances
    \item Radical arithmetic becomes circular or error-prone
    \item 3D visualization completely unclear
    \item Time exceeds 10 minutes without clear path to solution
\end{itemize}

\textbf{Before Abandoning}:
\begin{itemize}
    \item Try coordinate approach (more mechanical, less intuition needed)
    \item Check if answer can be guessed from answer choice patterns
    \item Verify dimensional analysis to eliminate some choices
\end{itemize}

\end{competitionbox}

\section*{Summary}

\subsection*{Key Takeaways}

\begin{enumerate}
    \item \textbf{Symmetry is Powerful}: The right pyramid structure immediately gives $|AV| = |DV|$ by symmetry, which is crucial for the solution.
    
    \item \textbf{45-45-90 Recognition}: Isosceles right triangles have the special property that legs are $\frac{\text{hypotenuse}}{\sqrt{2}}$. Recognizing this pattern saves significant time.
    
    \item \textbf{Regular Polygon Formulas}: For a regular $n$-gon, the central angle is $\frac{360^\circ}{n}$. Memorizing this helps with quick calculations.
    
    \item \textbf{3D Pythagorean Theorem}: In a right pyramid, $|AV|^2 = r^2 + h^2$ where $r$ is the base radius and $h$ is height. This is the bridge between 2D and 3D.
    
    \item \textbf{Multiple Methods Validate}: Both geometric reasoning and vector dot product give the same answer, increasing confidence.
    
    \item \textbf{Radical Simplification}: Recognizing patterns like $\sqrt{3+2\sqrt{2}} = 1 + \sqrt{2}$ (since $(1+\sqrt{2})^2 = 3 + 2\sqrt{2}$) is essential for late-band problems.
\end{enumerate}

\subsection*{Skills Practiced}

\begin{itemize}
    \item Regular polygon properties and calculations
    \item 3D spatial visualization
    \item Pythagorean theorem in three dimensions
    \item Law of Cosines
    \item Special right triangles (45-45-90)
    \item Vector dot product for perpendicularity
    \item Coordinate geometry in 3D
    \item Manipulation of expressions with radicals
    \item Symmetry exploitation
\end{itemize}

\subsection*{Final Answer}

The square of the height of the pyramid is:

\begin{center}
\fbox{\textbf{Answer: (B) $\displaystyle\frac{1+\sqrt{2}}{2}$}}
\end{center}

\textbf{Key Geometric Facts}:
\begin{itemize}
    \item Distance $|AD| = 1 + \sqrt{2}$ (chord of octagon)
    \item Circumradius $r^2 = \frac{2+\sqrt{2}}{2}$
    \item Triangle $AVD$ is 45-45-90 with legs $|AV| = |DV| = \frac{2+\sqrt{2}}{2}$
    \item Height: $h^2 = |AV|^2 - r^2 = \frac{1+\sqrt{2}}{2}$ \checkmark
\end{itemize}

\end{document}